%
% board_ordering_gate_only_buggy.tex
%
% FIDELITY CONTROL for docs/board_ordering.tex -- candidate fix (a) on its own.
%
% The applet tracks the highest BoardOrder place it has pushed and skips any
% push at a place no newer.  The counter is claimed under a lock; the socket
% write is not, because it cannot be.  This variant exists to show that
% claiming a place and writing the glass are two acts, and that ordering the
% first does not order the second.
%
% It differs from board_ordering_current_buggy.tex by one component (pushed)
% and one guard per staging operation.
%
% Type-check:  fuzz -t docs/board_ordering_gate_only_buggy.tex
% Model-check: probcli docs/board_ordering_gate_only_buggy.tex -model_check \
%                -p DEFAULT_SETSIZE 2 -p MAXINT 4 -goal "shown /= evershown"
%              (must report the goal FOUND)
%

\documentclass[a4paper,10pt,fleqn]{article}
\usepackage[margin=1in]{geometry}
\usepackage{fuzz}
\usepackage[colorlinks=true,linkcolor=blue,citecolor=blue,urlcolor=blue]{hyperref}

\begin{document}

\title{A Monotone Push Counter Alone: a Fidelity Control}
\author{Why claiming a place does not order the writes}
\date{August 2026}
\maketitle

% ============================================================
\section{Purpose}
% ============================================================

The first candidate fix has the applet remember the highest place it has
pushed and skip any push at a place no newer than that. The counter would sit
beside the slot and be claimed under a lock, in the same style as
\texttt{BoardSlot.store}.

It does not work, and the reason is structural rather than incidental. A push
is two acts: settling on what to write, and writing it. The counter orders the
first. The second is a socket round trip, which is exactly why it cannot be
done under a lock, so it stays unordered --- and it is the second act that the
user sees. Two pushers claim in the order $1, 2$ and land in the order $2, 1$,
and the glass comes to rest showing $1$ while the counter records $2$.

The state, \texttt{Init}, \texttt{TakePlace} and \texttt{StoreBoard} are those
of \texttt{docs/board\_ordering.tex}, plus one component.

% ============================================================
\section{Types and State}
% ============================================================

\begin{zed}
  [WORKER]
\end{zed}

\begin{zed}
  PLACE == 0 \upto 3
\end{zed}

\begin{schema}{Board}
  began : PLACE \\
  loading : WORKER \pfun PLACE \\
  held : PLACE \\
  shown : PLACE \\
  staged : WORKER \pfun PLACE \\
  pushed : PLACE \\
  everheld : PLACE \\
  evershown : PLACE \\
  offered : PLACE \\
  overtaken : \power PLACE
\where
  held \leq began \\
  shown \leq began \\
  pushed \leq began \\
  everheld \leq began \\
  evershown \leq began \\
  offered \leq began
\end{schema}

\texttt{pushed} is the candidate's counter: the highest place any push has
claimed. Every other component is as in \texttt{docs/board\_ordering.tex}.

\begin{schema}{Init}
  Board'
\where
  began' = 0 \\
  loading' = \emptyset \\
  held' = 0 \\
  shown' = 0 \\
  staged' = \emptyset \\
  pushed' = 0 \\
  everheld' = 0 \\
  evershown' = 0 \\
  offered' = 0 \\
  overtaken' = \emptyset
\end{schema}

\begin{schema}{TakePlace}
  \Delta Board \\
  w? : WORKER
\where
  w? \notin \dom loading \\
  began < 3 \\
  began' = began + 1 \\
  loading' = loading \cup \{ w? \mapsto began' \} \\
  held' = held \\
  shown' = shown \\
  staged' = staged \\
  pushed' = pushed \\
  everheld' = everheld \\
  evershown' = evershown \\
  offered' = offered \\
  overtaken' = overtaken
\end{schema}

\begin{schema}{StoreBoard}
  \Delta Board \\
  w? : WORKER
\where
  w? \in \dom loading \\
  ( loading~w? > held \land held' = loading~w? \\
  \quad~ \land overtaken' = overtaken ) \lor \\
  \quad~ ( loading~w? \leq held \land held' = held \\
  \quad~ \land overtaken' = overtaken \cup \{ loading~w? \} ) \\
  ( held' > everheld \land everheld' = held' ) \lor \\
  \quad~ ( held' \leq everheld \land everheld' = everheld ) \\
  loading' = \{ w? \} \ndres loading \\
  began' = began \\
  shown' = shown \\
  staged' = staged \\
  pushed' = pushed \\
  evershown' = evershown \\
  offered' = offered
\end{schema}

% ============================================================
\section{The Counter}
% ============================================================

Each staging operation now claims its place first and proceeds only if the
claim beats the counter. A claim that does not beat it disables the operation,
which is the skip.

\texttt{AnswerStage} claims the place of the board the slot holds.

\begin{schema}{AnswerStage}
  \Delta Board \\
  w? : WORKER
\where
  w? \notin \dom staged \\
  held > pushed \\
  pushed' = held \\
  staged' = staged \cup \{ w? \mapsto held \} \\
  ( held > offered \land offered' = held ) \lor \\
  \quad~ ( held \leq offered \land offered' = offered ) \\
  began' = began \\
  loading' = loading \\
  held' = held \\
  shown' = shown \\
  everheld' = everheld \\
  evershown' = evershown \\
  overtaken' = overtaken
\end{schema}

\texttt{RefreshStage} claims the place of the board this worker loaded.

\begin{schema}{RefreshStage}
  \Delta Board \\
  w? : WORKER
\where
  w? \in \dom loading \\
  w? \notin \dom staged \\
  loading~w? > pushed \\
  pushed' = loading~w? \\
  staged' = staged \cup \{ w? \mapsto loading~w? \} \\
  ( loading~w? > offered \land offered' = loading~w? ) \lor \\
  \quad~ ( loading~w? \leq offered \land offered' = offered ) \\
  began' = began \\
  loading' = loading \\
  held' = held \\
  shown' = shown \\
  everheld' = everheld \\
  evershown' = evershown \\
  overtaken' = overtaken
\end{schema}

\texttt{BlankStage} claims place $0$ --- the placeholder, or the message naming
why \texttt{bd} could not be read.

\begin{schema}{BlankStage}
  \Delta Board \\
  w? : WORKER
\where
  w? \notin \dom staged \\
  held = 0 \\
  0 > pushed \\
  pushed' = 0 \\
  staged' = staged \cup \{ w? \mapsto 0 \} \\
  began' = began \\
  loading' = loading \\
  held' = held \\
  shown' = shown \\
  everheld' = everheld \\
  evershown' = evershown \\
  offered' = offered \\
  overtaken' = overtaken
\end{schema}

\begin{schema}{PushCommit}
  \Delta Board \\
  w? : WORKER
\where
  w? \in \dom staged \\
  shown' = staged~w? \\
  ( staged~w? > evershown \land evershown' = staged~w? ) \lor \\
  \quad~ ( staged~w? \leq evershown \land evershown' = evershown ) \\
  staged' = \{ w? \} \ndres staged \\
  began' = began \\
  loading' = loading \\
  held' = held \\
  pushed' = pushed \\
  everheld' = everheld \\
  offered' = offered \\
  overtaken' = overtaken
\end{schema}

% ============================================================
\section{What it reaches}
% ============================================================

\begin{verbatim}
  probcli docs/board_ordering_gate_only_buggy.tex -model_check -nodead \
    -p DEFAULT_SETSIZE 2 -p MAXINT 4 -p TIME_OUT 60000 \
    -goal "shown /= evershown"                          # I2: FOUND
  probcli ... -goal "staged = {} & shown < offered"     # I3: FOUND
  probcli ... -goal "shown > held"                      # I4: FOUND
  probcli ... -goal "shown : overtaken"                 # I5: FOUND
  probcli ... -goal "held /= everheld"                  # I1: not found
\end{verbatim}

All four defects of the current design remain reachable. The counter closes
none of them. That is the result: it is not a partial fix to be completed
later, it is a field that costs a lock acquisition per push and changes nothing
a user can see.

The I2 witness is the two-act argument made concrete. Over 530 reachable states
ProB reaches it in eight steps: a click claims place $1$; a later load stores
place $3$ and a second click claims that; the second click's write lands, and
then the first click's write lands over it. The counter reads $3$ while the
display shows $1$. Claiming a place ordered the claims; the glass is written by
the other act.

\paragraph{The bounded-carrier deadlock.}
The bare \texttt{-model\_check} reports a deadlock, and it is an artifact of
the bound rather than a finding: once \texttt{pushed} reaches the top place, no
claim can beat it and no staging operation is enabled. Real places come from an
unbounded counter, so the state has no counterpart in the running system.
\texttt{-nodead} suppresses it so the goal checks run to exhaustion; 9 of the
531 states are its variants.

\paragraph{A second reason not to implement it.}
\texttt{BlankStage} is never covered. Its claim is $0 > pushed$, and
\texttt{pushed} is a place, so it is never below $0$. A strict monotone counter
over places would therefore suppress the placeholder and the failure message
outright --- including the very first placeholder of a session, the one a cold
click depends on. Any counter of this shape needs a value strictly below the
bottom place to start from, which is one more piece of machinery justifying a
mechanism that does not work anyway. ProB reports the operation not covered,
which is the cheapest possible way to be told.

\end{document}
